\documentclass[12pt,a4paper]{article}
\usepackage[utf8]{inputenc}
\usepackage{geometry}
\geometry{left=1in, right=1in, top=1in, bottom=1in}
\usepackage{hyperref}
\usepackage{titlesec}
\usepackage{enumitem}

% Formatting hyper-links to look clean in TOC
\hypersetup{
    colorlinks=true,
    linkcolor=black,
    filecolor=black,
    urlcolor=black,
}

\begin{document}

% --- TITLE PAGE ---
\begin{titlepage}
    \centering
    \vspace*{1.5cm}
    {\Huge \textbf{Home Assignment 1}\par}
    \vspace{1.5cm}
    {\Large Ten Steps in Case Analysis for [Insert Case Topic Here]\par}
    \vspace{2cm}
    {\large ENGR4960J\par}
    \vspace{0.5cm}
    {\large Date: [Insert Date]\par}

    \vfill

    {\textbf{Group [\#]: [Name 1], [Name 2], [Name 3], [Name 4], [Name 5]}\par}
    \vspace{0.5cm}

    {\textbf{Contribution of each group member}\par}
    \begin{flushleft}
    \textbf{[Name 1]:} [Contribution description...]\par
    \textbf{[Name 2]:} [Contribution description...]\par
    \textbf{[Name 3]:} [Contribution description...]\par
    \textbf{[Name 4]:} [Contribution description...]\par
    \textbf{[Name 5]:} [Contribution description...]\par
    \end{flushleft}

    \vspace{0.5cm}

    {\textbf{AI statement:}\par}
    \begin{flushleft}
    [AI statement (if AI was used within the workflow – you can't use AI to generate the text. Your final text will be checked through AI detector)]
    \end{flushleft}

    \vspace{1cm}
\end{titlepage}

% --- TABLE OF CONTENTS ---
\tableofcontents
\newpage

% --- General Case Description ---
\section*{Case introduction}
\addcontentsline{toc}{section}{Case introduction}

\vspace{0.3cm}
\noindent\textbf{General case description (up to 100 words):} In this section, you can explain the context of the case and its general chronology. Try to provide a clear case description, as it can significantly improve your case analysis.

\vspace{0.5cm}

% --- STEP 1 ---
\section{Identifying Ethical Issues}
\vspace{0.3cm}
\noindent Describe why this step is important 

\subsection{Company-Related Issues}
\noindent\textbf{Issue 1: [Title of the issue]}\\
\textbf{Description of the issue:} [Description...]

\vspace{0.3cm}
\noindent\textbf{Issue 2, 3, n, ...}\\

\subsection{Engineers-Related Issues}
\noindent\textbf{Issue 1: [Title of the issue]}\\
\textbf{Description of the issue:} [Description...]

\vspace{0.3cm}
\noindent\textbf{Issue 2, 3, n, ...}\\

\subsection{Other Issues}
\noindent\textbf{Issue 1: [Title of the issue]}\\
\textbf{Description of the issue:} [Description...]

\vspace{0.3cm}
\noindent\textbf{Issue 2, 3, n, ...}\\

\vspace{0.5cm}
\noindent\textbf{Additional question: Who or what is being affected, and how?}\\
Here, list at least three distinct stakeholder groups and explain how each was affected by the accident. For example, possible stakeholder groups include employees, suppliers, customers, governments, local communities, and others.

% --- STEP 2 ---
\section{Narrowing the Focus}
\vspace{0.3cm}
\noindent Describe why this step is important 
\vspace{0.3cm}

\noindent We decided to narrow the focus to Issue \#[mention the issue here] due to the following reasons.

\vspace{0.3cm}
\noindent\textbf{This issue is a direct/indirect cause of the event because...} [Description...]

\vspace{0.3cm}
\noindent\textbf{The severity of the consequences}\\
\textbf{Description:} [Description...]

\vspace{0.3cm}
\noindent\textbf{The scope of the consequences}\\
\textbf{Description:} [Description...]

\vspace{0.3cm}
\noindent\textbf{Long-term and Short-term consequences}\\
\textbf{Description:} [Description...]

\vspace{0.3cm}
\noindent\textbf{Other reasons}\\
\textbf{Description:} [Description...]

% --- STEP 3 ---
\section{Determining Relevant and Missing Facts}
\vspace{0.3cm}
\noindent Describe why this step is important 

\vspace{0.3cm}
\noindent We categorized facts in accordance with the following types:

\subsection{General Facts}
\noindent\textbf{Date: } [Insert Date]\\
\textbf{Location: } [Insert Location]\\
\textbf{Other relevant facts: } [Insert Other Facts]

\subsection{Material Facts}
\textbf{Description:} [Description...]

\subsection{Facts Regarding Individuals}
\textbf{Description:} [Description...]

\subsection{Facts Regarding Organizations}
\textbf{Description:} [Description...]

\subsection{Missing Facts}
\noindent\textbf{Fact 1, 2, n, ...}\\
\textbf{Description:} [Description...]

% --- STEP 4 ---
\section{Reasonable Assumptions}
In real-life engineering practice, missing information is a common challenge. It is essential to make reasonable assumptions based on available facts, professional experience, and logical inference to make informed, ethical, and safety-focused decisions.

Relying on the missing information from the case (e.g., missing facts), we asked the following question to make a reasonable assumption.

\vspace{0.3cm}
\noindent\textbf{Question 1: } [Insert Question]\\
\textbf{Reasonable assumption 1: } [Insert Assumption]\\
\textbf{Description (reason for making an assumption): } [Description...]

\vspace{0.3cm}
\noindent\textbf{Question 2, 3, n, ...}\\
\textbf{Reasonable assumption 2, 3, n, ...}\\

% --- STEP 5 ---
\section{Undertaking Definitional Clarification}
\vspace{0.3cm}
\noindent Describe why this step is important (50-100 words)

\vspace{0.3cm}
\noindent (Include at least 3 terms)

\noindent Here is the list of the definitions we have to clarify:

\vspace{0.3cm}
\noindent\textbf{Term 1: } [Insert Term]\\
\textbf{Definition: } [Insert Definition]\\
\textbf{Clarification: } [Insert Clarification]\\
\textbf{Why is it relevant to the case?} [Insert Reason]

\vspace{0.3cm}
\noindent\textbf{Term 2, 3, n, ...}\\
\textbf{Definition: } [Insert Definition]\\
\textbf{Clarification: } [Insert Clarification]\\
\textbf{Why is it relevant to the case?} [Insert Reason]

% --- STEP 6 ---
\section{Step 6. Applying Ethical Principles}
Please refer to the list of principles for global engineering. List as many as possible.

\vspace{0.3cm}
\noindent Describe why this step is important (50-100 words).

\vspace{0.3cm}
\noindent\textbf{Principle \#... Title of the principle} (for example, ‘Keep the public safe from the negative consequences...’)\\
\textit{Describe why you think it is applicable to the case :} 

\vspace{0.3cm}
\noindent\textbf{Principle 2, 3, n, ...} \\
\textit{Describe why you think it is applicable to the case :} 

\vspace{0.3cm}
\noindent\textbf{*Rationale for exclusion*:} if you didn’t include some principles from the list mentioned in a footnote, you should explain why in this sub-section.

\subsection*{Conflicts Between Principles}
Some principles may not agree with each other. For example, Principle 6 (honesty) and business interests can conflict. Find and list possible conflicts between the principles and explain how to solve them from the perspective of engineering ethics.

% --- STEP 7 ---
\section{Step 7. Review process}
(list as many as possible)

\vspace{0.3cm}
\noindent Describe why this step is important (50-100 words)

\vspace{0.3cm}
\noindent \textbf{During the review process, we found several essential elements that hadn’t been mentioned during the analysis procedure earlier :}

\vspace{0.3cm}
\noindent \textbf{Element 1}, Description

\vspace{0.3cm}
\noindent \textbf{Element 2}, Description

\vspace{0.3cm}
\noindent \textbf{Element 3}, Description

% --- STEP 8 ---
\section{Step 8. Solving Ethical Issues}
(Solve as many as possible)

\vspace{0.3cm}
\noindent Describe why this step is important (50-100 words)

\vspace{0.3cm}
\noindent In this step, you must identify several issues you mentioned in the previous step (step 1) and provide your solution to them.

\vspace{0.3cm}
\noindent \textbf{Within the first step, we have identified the following issues:}

\vspace{0.3cm}
\noindent \textbf{Issue 1, Title;}\\
Short-term solution -- \\
Long-term solution -- 

\vspace{0.3cm}
\noindent \textbf{Issue 2, Title; Solution --}\\
Short-term solution -- \\
Long-term solution -- 

\vspace{0.3cm}
\noindent \textbf{Issue 3, Title; Solution --}\\
Short-term solution -- \\
Long-term solution -- 

\vspace{0.5cm}
\noindent \textbf{There are some other issues that we haven’t identified at the very beginning, but the importance of these issues became evident after the review process. Here is the list:}

\vspace{0.3cm}
\noindent \textbf{Issue 1, Title;}\\
Short-term solution -- \\
Long-term solution -- 

\vspace{0.3cm}
\noindent \textbf{Issue 2, Title;}\\
Short-term solution -- \\
Long-term solution -- 

\vspace{0.3cm}
\noindent \textbf{Issue 3, Title; Solution --}\\
Short-term solution -- \\
Long-term solution -- 

% --- STEP 9 ---
\section{Step 9. Identifying practical constraints}
(list as many as possible)

\vspace{0.3cm}
\noindent Describe why this step is important (50-100 words)

\subsection*{Practical constraints related to technological \& engineering part of the case}
\noindent \textbf{Practical constrain 1; Name of the constrain;} Description (+the importance for this case)\\
\noindent \textbf{Practical constrain 2; Name of the constrain;} Description (+the importance for this case)\\
\noindent \textbf{Practical constrain 3; Name of the constrain;} Description (+the importance for this case)

\subsection*{Practical constraints related to Regulatory and Legal Frameworks}
\noindent \textbf{Practical constrain 1; Name of the constrain;} Description (+the importance for this case)\\
\noindent \textbf{Practical constrain 2; Name of the constrain;} Description (+the importance for this case)\\
\noindent \textbf{Practical constrain 3; Name of the constrain;} Description (+the importance for this case)

\subsection*{Practical constraints related to the cultural environment}
\noindent \textbf{Practical constrain 1; Name of the constrain;} Description (+the importance for this case)\\
\noindent \textbf{Practical constrain 2; Name of the constrain;} Description (+the importance for this case)\\
\noindent \textbf{Practical constrain 3; Name of the constrain;} Description (+the importance for this case)

\subsection*{Other practical constraints}
\noindent \textbf{Practical constrain 1; Name of the constrain;} Description (+the importance for this case)\\
\noindent \textbf{Practical constrain 2; Name of the constrain;} Description (+the importance for this case)\\
\noindent \textbf{Practical constrain 3; Name of the constrain;} Description (+the importance for this case)

% --- STEP 10 ---
\section{Step 10. Avoiding/Predicting Ethical Issues in the future}
Use your moral imagination to list possible ethical issues that this technology could raise in the future. Please include problems that we haven’t explicitly discussed before. Also, provide engineering and design solutions that might help prevent this problem from happening.

\vspace{0.3cm}
\noindent Describe why this step is important (50-100 words)

\vspace{0.3cm}
\noindent \textbf{Issue 1; Description;} Engineering Solution + Governmental/Legal solution + other possible solutions

\vspace{0.3cm}
\noindent \textbf{Issue 2; Description;} Engineering Solution + Governmental/Legal solution + other possible solutions

\vspace{0.3cm}
\noindent \textbf{Issue 3; Description;} Engineering Solution + Governmental/Legal solution + other possible solutions



% --- CONCLUSION ---
\section{Self-Reflection report }

\textit{Question \#1}:\textbf{ }“How did the AI participate in deciding what your artifact became?” 

\vspace{0.3cm}
\noindent \textbf{Answer (150 words):}

\vspace{0.3cm}
\noindent \textit{Question \#2:} What conversational patterns did you find particularly useful while designing your program with AI? How could these conversational patterns be different from your conversations with humans? 

\vspace{0.3cm}
\noindent \textbf{Answer (150 words)}:

\vspace{0.3cm}
\noindent \textit{Question \#3}: Find one place in your tool's code or copy where you accepted something the AI generated without checking whether it was accurate or actually what you meant. What is it? 

\vspace{0.3cm}
\noindent \textbf{Answer (150 words):}


% --- REFERENCES ---
\newpage
\section*{References}
\addcontentsline{toc}{section}{References}
You must add references to your work (at least 10). Here are some requirements:
\begin{itemize}[itemsep=0pt, parsep=0pt]
    \item References should be situated on a separate page.
    \item References in APA 7 – I recommend using the Mendeley citation tool for convenience and accuracy.
    \item The list of references should be situated in alphabetic order.
    \item References should be provided to the scientific literature (like articles, monographs, etc).
    \item Please do not refer to Wikipedia or other sources that are doubtful.
    \item You also cannot make references on ChatGPT.
    \item You can refer to both English and Chinese references.
    \item All the references mentioned in the list should be quoted in the paper.
\end{itemize}

\vspace{0.5cm}
\noindent \textbf{Appendix 1}\\
\textbf{Example for the list of references in APA7}

\vspace{0.3cm}
\noindent \textit{Please note that the references included in this list are merely examples and do not directly relate to your case analysis. You should include your references.}

\vspace{0.5cm}
\noindent 1. Coeckelbergh, M. (2017). Can Machines Create Art?. \textit{Philos. Technol. 30}, 285–303 https://doi.org/10.1007/s13347-016-0231-5

\vspace{0.3cm}
\noindent 2. Coeckelbergh, M. (2023). The Work of Art in the Age of AI Image Generation: Aesthetics and Human-Technology Relations as Process and Performance. \textit{Journal of Human-Technology Relations, 1}. https://doi.org/10.59490/jhtr.2023.1.7025

\vspace{0.3cm}
\noindent 3. Messeri, L., Crockett, M.J. (2024) Artificial intelligence and illusions of understanding in scientific research. \textit{Nature 627}, 49–58 https://doi.org/10.1038/s41586-024-07146-0

\vspace{0.3cm}
\noindent 4. Wang, H., Fu, T., Du, Y. et al. (2023) Scientific discovery in the age of artificial intelligence. \textit{Nature 620}, 47–60 https://doi.org/10.1038/s41586-023-06221-2

\end{document}
